% ----- CONFIGURAÇÕES ----- %
\documentclass[12pt]{article}
\usepackage{graphicx} % imagens
\usepackage{lmodern} % versão vetorial das fontes internas do LaTeX (evita "Font shape not available" em tamanhos customizados como o título da capa)
\usepackage[brazilian]{babel} % idioma
\usepackage{pdflscape}
\usepackage{amsmath} % equações
\numberwithin{equation}{section} % numeração de equação numero-da-secao.numero-da-equacao
\usepackage{caption} % legenda em imagens e tabelas
\captionsetup{justification=centering} % legendas e fontes sempre centralizadas, com uma ou várias linhas
\usepackage{subcaption} % legenda em subimagens e subtabelas
\newcommand{\source}[1]{\caption*{Fonte: {#1}} } % para adicionar fonte de imagens e tabelas
\usepackage{enumitem} % criação de listas
\usepackage{cite} % para citar seções, imagens...
\usepackage[hidelinks]{hyperref} % para ter itens (citação de seções, sumário...) clicáveis
\usepackage{indentfirst} % espaçamento do primeiro parágrafo
\setlength{\parindent}{1.25cm} % recuo de primeira linha de parágrafo (ABNT NBR 14724)
\setcounter{section}{-1} % começa as seções em 0
\usepackage{booktabs}
% --- ABNT --- %
% definindo fonte
\usepackage{fontspec}
\setmainfont{Times New Roman}
% tamanho da folha e margens
\usepackage[a4paper, left=3cm, top=3cm, right=2cm, bottom=2cm, headheight=15pt]{geometry}
% definindo tamanho para items e subitems
\usepackage{titlesec}
\titleformat{\section}{\normalfont\fontsize{16}{19.2}\bfseries}{\makebox[25pt][l]{\thesection}}{0pt}{}
\titleformat{\subsection}{\normalfont\fontsize{16}{19.2}\bfseries}{\makebox[25pt][l]{\thesubsection}}{0pt}{}
% para configurar alinhamento e espaçamento
\usepackage{setspace}
\usepackage{ragged2e}
\titlespacing*{\section}{0cm}{12pt}{12pt}
\titlespacing*{\subsection}{0cm}{12pt}{12pt}
\usepackage[alf,abnt-emphasize=bf]{abntex2cite} % bibliografia e citação no padrão ABNT
% --- numeração de página (ABNT NBR 14724): conta desde a capa, mas só aparece a partir da introdução --- %
\usepackage{fancyhdr}
\pagestyle{fancy}
\fancyhf{}
\fancyhead[R]{\thepage}
\renewcommand{\headrulewidth}{0pt}
\pagestyle{empty} % elementos pré-textuais (capa, resumo, abstract, listas, sumário) não exibem número
% --- pacotes extras: tabelas, listas, texto e cores mais avançados (ver README.md, seção 6) --- %
\usepackage{array}      % colunas de tabela com largura fixa e quebra de linha (p{}, m{})
\usepackage{tabularx}   % colunas de tabela que se ajustam automaticamente à largura da página
\usepackage{multirow}   % mesclar células verticalmente em tabelas
\usepackage{longtable}  % tabelas que quebram automaticamente entre páginas
\usepackage{xcolor}     % texto colorido / destacado
\usepackage[normalem]{ulem} % texto sublinhado (\uline) e tachado (\sout)
\usepackage{wrapfig}    % texto correndo ao redor de uma imagem
\usepackage{amssymb}    % símbolos matemáticos extras
\usepackage{listings}   % blocos de código-fonte formatados
\usepackage{float}
\usepackage{pgfplots} % graficos desenhados no proprio LaTeX (Figura de receita na Secao de mercado)
\pgfplotsset{compat=1.18}
\usetikzlibrary{babel} % evita conflito entre TikZ e os atalhos do babel em portugues
\usepackage[section]{placeins} % figuras e tabelas não atravessam o início de uma nova seção; \FloatBarrier fixa o limite dentro dela

\lstset{ % estilo padrão dos blocos de código (ver README.md, seção 6)
    basicstyle=\ttfamily\small,
    backgroundcolor=\color{gray!10},
    frame=single,
    breaklines=true,
    numbers=left,
    numberstyle=\tiny\color{gray},
    showstringspaces=false
}
